\chapter{Аудит кандидатов odd auxiliary cross-bilinear}
\label{ch:v10-h15-r2-odd-auxiliary-cross-bilinear-audit}

Hubbard--Stratonovich identity свёл оставшуюся проблему к одному оператору:
требуется независимо вывести mixed bilinear
\begin{equation}
 \langle A_\Sigma,Q\Sigma\rangle,
 \qquad Q=6Y.
\end{equation}
На восьми секторах
\begin{equation}
 T=6T_R^3=(3,-3,3,-3,3,-3,3,-3),
 \quad B=3(B-L)=(0,0,4,4,-4,-4,0,0),
\end{equation}
и $Q=T+B=(3,-3,7,1,-1,-7,3,-3)$. Два Cartan-столбца независимы,
их Gram matrix равна $\operatorname{diag}(72,64)$, а единственная
комбинация, дающая $Q$, имеет коэффициенты $(1,1)$.

Двенадцать механизмов проверены по шести критериям: правильный
$A_\Sigma$--$\Sigma$ domain, точный $Q$, gauge equivariance, Real-even
scalar, фиксация коэффициента одним parent и присутствие в унаследованном
action. Матрица имеет rank $6$, но pass-vector равен $0_{12}$.

Universal portal даёт только $I$. Fixed-point projection и старый
mapping-cone incidence унаследованы, но не имеют нужного domain и $Q$.
Ordinary even spectral moments не создают линейный Cartan operator.
Коммутатор $[D,Y]$ исчезает на гиперзаряд-сохраняющем seed. Два отдельных
Cartan portal воспроизводят $Q$ лишь после ручной установки отношения
$1:1$. KO6 fluctuation присутствует, но его текущая опора не содержит
полного $Q$-канала. Callias и fermion-loop варианты оставляют свободную
нормировку.

Два кандидата получают $5/6$. Первый --- trilinear, построенный одним
moment map нейтрального $\Delta$-фона,
\begin{equation}
 \langle A_\Sigma,(6\mu_R-4\operatorname{ad}\mu_4)\Sigma\rangle
 =\langle A_\Sigma,Q\Sigma\rangle.
\end{equation}
Второй --- соответствующая mixed component superconnection curvature.
Оба автоматически фиксируют Cartan-комбинацию, но ни один не входит в
текущий mapping-cone action вместе с независимым $A_\Sigma$.

\begin{theorem}[Узкий источник cross-оператора]
В Cartan-span оператор $Q$ единственен и требует коэффициенты $(1,1)$.
Из двенадцати проверенных механизмов ни один не проходит inherited-action
criterion. Ближайший минимальный путь --- единый $\Delta$ moment-map
trilinear, эквивалентно mixed superconnection curvature, с оценкой $5/6$.
\end{theorem}

\begin{nogocriterion}[Два Cartan portal не являются одним parent]
Нельзя считать $Q=T+B$ выведенным, если коэффициенты двух независимых
bilinear устанавливаются после просмотра target spectrum. Отношение должно
следовать из одного нормированного moment map до Schur elimination.
\end{nogocriterion}

\begin{protocol}\begin{enumerate}
\item Cartan basis rank: \textbf{$2$};
\item locked coefficients $(T,B)$: \textbf{$(1,1)$};
\item inherited/required cross rank: \textbf{$0/8$};
\item full cross-Hessian rank: \textbf{$16$};
\item cross-Hessian inertia: \textbf{$(8_-,0_0,8_+)$};
\item audit rank: \textbf{$6$};
\item прошедших кандидатов: \textbf{$0/12$};
\item лучшие кандидаты: \textbf{$\Delta$ moment map и superconnection, $5/6$};
\item physical origin: \textbf{$5/6$};
\item следующий гейт: \textbf{родитель $\Delta$ moment-map auxiliary trilinear}.
\end{enumerate}\end{protocol}

Машинный сертификат сохранён в
\path{s2t_v10_particle_wrinkle_dislocation_callias_equal_charge_twist_m4_h15_r2_odd_auxiliary_cross_bilinear_candidate_audit_gate_results.json}.